\documentclass[11pt,a4paper]{article}
\usepackage[utf8]{inputenc}
\usepackage[T1]{fontenc}
\usepackage[brazil]{babel}
\usepackage[margin=2.5cm]{geometry}
\usepackage{mathptmx}
\usepackage{enumitem}
\usepackage{booktabs}
\usepackage{tabularx}
\usepackage{array}
\usepackage{hyperref}
\hypersetup{colorlinks=false, pdfborder={0 0 0}}
\usepackage{fancyhdr}

\pagestyle{fancy}
\fancyhf{}
\fancyhead[L]{\small TrackFleet --- Plano de Engajamento}
\fancyhead[R]{\small\thepage}
\renewcommand{\headrulewidth}{0.4pt}

\newcolumntype{L}{>{\raggedright\arraybackslash}X}
\newcolumntype{C}{>{\centering\arraybackslash}X}

\begin{document}

\begin{center}
{\Large\bfseries DOCUMENTO DE PARTES INTERESSADAS DO PROJETO}\\[2pt]
{\large\bfseries TrackFleet --- Gestão de Rastreador de Automóveis}\\[4pt]
Disciplina de Gerenciamento de Projetos $\cdot$ Framework PMBOK\textregistered\ 8\textsuperscript{a} Edição $\cdot$ Domínio: Partes Interessadas\\[2pt]
Integrantes: Enzo Allebrand e Kerlon Ribeiro (dois GPs) $\cdot$ 4\textsuperscript{o} semestre $\cdot$ Data: 27/08
\end{center}

\vspace{2mm}
\hrule
\vspace{4mm}

\noindent\textbf{\large Plano de Engajamento}

\vspace{1mm}
\noindent Define o nível atual e desejado de engajamento de cada parte interessada, estabelecendo as estratégias e ações para fechar eventuais lacunas e garantir o suporte necessário ao projeto.

\section{Matriz de Avaliação do Engajamento}

Níveis de engajamento: \textbf{D} (Desinformado), \textbf{R} (Resistente), \textbf{N} (Neutro), \textbf{A} (Apoiador), \textbf{L} (Líder).
Legenda: \textbf{A} = Nível Atual, \textbf{D} = Nível Desejado.

\begin{table}[h]
\centering
\small
\begin{tabularx}{\textwidth}{@{} L C C C C C @{}}
\toprule
\textbf{Parte Interessada} & \textbf{Desinformado} & \textbf{Resistente} & \textbf{Neutro} & \textbf{Apoiador} & \textbf{Líder} \\
\midrule
Diretoria da Transportadora & & & & A & D \\
GPs / Equipe (Enzo e Kerlon) & & & & & A, D \\
Gestores de Frota & A & & & D & \\
Motoristas & & A & D & & \\
Fornecedores de APIs/Hardware & & & A, D & & \\
\bottomrule
\end{tabularx}
\end{table}

\section{Estratégias de Engajamento e Ações}

\begin{itemize}[itemsep=4pt]
    \item \textbf{Diretoria (Apoiador $\rightarrow$ Líder):}
    \begin{itemize}[label=--]
        \item \textbf{Ação:} Promover reuniões curtas e executivas para validar resultados de negócio.
        \item \textbf{Responsável:} GPs.
        \item \textbf{Gatilho de reavaliação:} Caso o patrocinador pare de responder solicitações ou falte às reuniões de aprovação.
    \end{itemize}
    
    \item \textbf{Gestores de Frota (Desinformado $\rightarrow$ Apoiador):}
    \begin{itemize}[label=--]
        \item \textbf{Ação:} Envolvê-los ativamente na definição e testes dos relatórios de alertas. Demonstrações dos incrementos de software.
        \item \textbf{Responsável:} GPs.
        \item \textbf{Gatilho de reavaliação:} Baixa adesão nas fases de testes de aceitação ou feedbacks genéricos.
    \end{itemize}

    \item \textbf{Motoristas (Resistente $\rightarrow$ Neutro):}
    \begin{itemize}[label=--]
        \item \textbf{Ação:} Garantir a comunicação transparente de que o sistema visa segurança (redução de furtos/acidentes) e não punição. Criar termos de uso e privacidade transparentes (conformidade LGPD).
        \item \textbf{Responsável:} Patrocinador, com apoio técnico dos GPs.
        \item \textbf{Gatilho de reavaliação:} Reclamações trabalhistas ou sabotagem aos equipamentos rastreadores.
    \end{itemize}
\end{itemize}

\end{document}
